%----------------------------------------------------------------------------------------
%	                            PORTADA
%----------------------------------------------------------------------------------------
% \makeportada{}{Daniel Carbajales Soto}{ 1° CFGS Automatización y Robótica Industrial}{}

%----------------------------------------------------------------------------------------
%                                    HEADER 
%----------------------------------------------------------------------------------------
% \header{}{Daniel Carbajales Soto}{1° CFGSARI}
%----------------------------------------------------------------------------------------
%                               TABLE OF CONTENTS
%----------------------------------------------------------------------------------------
% \maketoc
%----------------------------------------------------------------------------------------
%	                            TABLES
%----------------------------------------------------------------------------------------
% this is just a demo in case i need to make a table
%
%
%
% \begin{tabularx}{\linewidth}{|C{2cm}|Y|Y|C{2cm}|}
% \hline
% \makecell{Header 1} & \makecell{Header 2 \\ with linebreak} & \makecell{Header 3} & \makecell{Header 4} \\
% \hline
% Row 1 & This is a long cell text that automatically wraps in the Y column & Another long cell & Short \\
% \hline
% Row 2 & \multirow{2}{*}{Multirow text} & More text & Short \\
% \cline{1-1} \cline{3-4}
% Row 3 & & Extra text & Short \\
% \hline
% \end{tabularx}


%----------------------------------------------------------------------------------------
%	                            DOCUMENT
%----------------------------------------------------------------------------------------
\fchapter{3}{Pág. 23}
\fsection{\boldit{\textcolor{waveBlue2}{Actividades:} Ponte a prueba}}
\addcontentsline{toc}{section}{\textcolor{waveBlue2}{Ejercicio 7}}
\section*{\textcolor{waveBlue2}{7. }Encuentra los tres errores que aparecen en cada uno de los siguientes textos y redáctalos de manera correcta.}

\subsection*{\boldit{\textcolor{waveBlue1}{Texto 1:}}}
La transformación digital y la  digitalización son conceptos distintos pero relacionados. Mientras que la  \textcolor{samuraiRed}{transformación digital} $\to$ ``digitalización'' se 
enfoca en convertir procesos y datos a formatos digitales para modernizar operaciones, la \textcolor{samuraiRed}{digitalización} $\to$ ``transformación digital'' implica un cambio profundo en el uso de
tecnología para mejorar el rendimiento y lograr objetivos estratégicos, impactando modelos de negocio, cultura organizacional y procesos. 
Esta transformación puede manifestarse en la reingeniería de técnicas empresariales, la evolución del modelo de negocio, 
la adaptación al dominio empresarial y el cambio cultural organizacional, \textcolor{samuraiRed}{del mismo modo a todas las empresas} $\to$ ``adaptándose personalizadamente a cada una de las empresas.''

\subsection*{\boldit{\textcolor{waveBlue1}{Texto 2:}}}
La transformación \textcolor{samuraiRed}{digital parcial} $\to$ ``digital '' de una empresa abarca la aplicación de tecnologías digitales en todos los aspectos del negocio, 
desde la producción hasta las operaciones empresariales y la \textcolor{samuraiRed}{cadena alimentaria} $\to$ ``cadena de suministro''. Esta \textcolor{samuraiRed}{digitalización} $\to$``transformación'' 
mejora la eficiencia operativa mediante la documentación de procesos y la monitorización en tiempo real, con tecnologías como IoT y sistemas de control avanzado. 
El proceso de transformación digital tiene asociado un desafío en su implantación llamado resistencia al cambio, que se supera 
aplicando una comunicación efectiva, educando y haciendo partícipe a todo el personal.

%%%%%%%%%%%%%%%%%%%%%%%%%%%%%%%%%%%%%%%%%%%%%%%%%%%%%%%%%%
%%%%%%%%%%%%%%%%%%%% bibliography %%%%%%%%%%%%%%%%%%%%%%%%
% \nocite{*}                                 %%%%%%%%%%%%%
% \printbibliography                         %%%%%%%%%%%%%
%%%%%%%%%%%%%%%%%%%%%%%%%%%%%%%%%%%%%%%%%%%%%%%%%%%%%%%%%%
%%%%%% Last page in case there is no bibliography %%%%%%%%
% \label{Last Page}                           %%%%%%%%%%%%%
%%%%%%%%%%%%%%%%%%%%%%%%%%%%%%%%%%%%%%%%%%%%%%%%%%%%%%%%%%

